\documentclass[11pt,letterpaper]{article}
\usepackage{fontspec}
\usepackage{tocloft}
\usepackage{multicol}
\usepackage{nicefrac}
\usepackage[left=1.4in,right=1.5in,top=1in,bottom=1in]{geometry}
\setmainfont[Scale=1.3,AutoFakeBold=1.5,AutoFakeSlant=0.3]{Doves Type}

\widowpenalty=10000
\clubpenalty=10000
\interlinepenalty=500

\title{Raspberry Vinaigrette}
\author{}
\date{}

\begin{document}

\maketitle
\thispagestyle{empty}

\section*{Ingredients}

\setlength{\columnsep}{20pt}
\begin{multicols}{2}
\noindent
    Shallot \dotfill 1 \\
    Raspberry shrub \dotfill 5 Tbsp. \\
    Dijon mustard \dotfill 1 tsp. \\
    Capers \dotfill 1 tsp. \\
    \columnbreak
    Caper brine \dotfill 1 tsp. \\
    Extra-virgin olive oil \dotfill 5 Tbsp. \\
    Neutral oil \dotfill 4 Tbsp. \\
    Fresh dill \dotfill 2 Tbsp. \\
    Kosher salt \dotfill to taste \\
    Black pepper \dotfill to taste
\end{multicols}

\section*{Directions}

\noindent
Mince \textbf{shallot} and combine with \textbf{raspberry shrub} and \textbf{Dijon} in \textit{Small Bowl~\#1} --- Chop \textbf{capers} and measure \textbf{caper brine} in \textit{Small Bowl~\#2} --- Measure \textbf{olive oil} and \textbf{neutral oil} together in \textit{Medium Bowl~\#1} --- Chop \textbf{dill} and set aside in \textit{Small Bowl~\#3}

\begin{enumerate}
    \item In a medium mixing bowl, combine \textbf{shallot}, \textbf{raspberry shrub}, and \textbf{Dijon} (\textit{Small Bowl~\#1}) with \textbf{capers} and \textbf{caper brine} (\textit{Small Bowl~\#2}). Whisk for \textit{20--30~seconds} until the mustard fully dissolves and the mixture looks smooth and lightly thickened, with no visible mustard clumps.

    \item Let the base stand for \textit{5~minutes} at room temperature so the \textbf{shallot} softens slightly and loses raw harshness. The mixture is ready when the shallot smell is milder and more aromatic than sharp.

    \item Whisk constantly while adding \textbf{olive oil} and \textbf{neutral oil} (\textit{Medium Bowl~\#1}) in a thin stream over \textit{45--60~seconds}. Continue whisking another \textit{15--20~seconds} until the vinaigrette is glossy, lightly thickened, and holds together for a few seconds before slowly separating.

    \item Stir in \textbf{dill} (\textit{Small Bowl~\#3}) and season with \textbf{kosher salt} and \textbf{black pepper}. Taste and adjust: if it is too sweet, add 1~tsp. \textbf{white wine vinegar}; for a stronger raspberry profile, add 1~tsp. \textbf{raspberry shrub}; for more savory lift, add a few drops of \textbf{caper brine}; if it tastes too sharp, whisk in 1~tsp. \textbf{neutral oil}.

    \item Let the dressing rest for \textit{10~minutes} before serving so flavors integrate. Whisk briefly just before using; it should taste fruit-forward and lightly sweet up front, then finish herbal, briny, and clean.
\end{enumerate}

\newpage

\section*{Equipment Required}
\begin{itemize}
    \item Medium mixing bowl
    \item Whisk
    \item Measuring spoons and measuring cup
    \item Chef's knife and cutting board
    \item Three small prep bowls and one medium prep bowl
    \item Jar with lid for storage (optional)
\end{itemize}

\section*{Hints and Notes}

\subsection*{Yield}
\begin{itemize}
    \item Makes about 1~cup vinaigrette
    \item Serves about 8 (2~Tbsp. per serving)
\end{itemize}

\subsection*{Mise en Place}
\begin{itemize}
    \item \textit{Small Bowl~\#1}: minced \textbf{shallot}, \textbf{raspberry shrub}, and \textbf{Dijon}
    \item \textit{Small Bowl~\#2}: chopped \textbf{capers} and \textbf{caper brine}
    \item \textit{Medium Bowl~\#1}: \textbf{olive oil} and \textbf{neutral oil}
    \item \textit{Small Bowl~\#3}: chopped fresh \textbf{dill}
    \item Have \textbf{salt} and \textbf{black pepper} ready for final seasoning
\end{itemize}

\subsection*{Ingredient Tips}
\begin{itemize}
    \item Since the \textbf{shrub} is vinegar-based, it serves as both fruit and acid; add 1--2~tsp. \textbf{white wine vinegar} only if the final dressing tastes too sweet
    \item \textbf{White wine vinegar} is the best adjustment acid here; \textbf{red wine vinegar} is bolder and can overshadow dill
    \item Use a mix of \textbf{olive oil} and \textbf{neutral oil} so raspberry stays prominent while keeping good body
    \item Salt-packed or brined \textbf{capers} both work; if salt-packed, rinse and soak briefly before chopping
    \item Fresh \textbf{dill} is essential here; dried dill will taste flat and dusty
\end{itemize}

\subsection*{Preparation Tips}
\begin{itemize}
    \item Whisk the acid base first to dissolve \textbf{Dijon} completely before adding oil
    \item Add oil in a steady thin stream for a more stable emulsion and silkier mouthfeel
    \item Resting the dressing for \textit{10~minutes} noticeably rounds acidity and integrates dill and caper aromas
    \item Taste over a leaf of the salad greens you plan to use; bitterness changes perceived sweetness and salt
\end{itemize}

\subsection*{Make Ahead \& Storage}
\begin{itemize}
    \item Best texture and fresh dill aroma are at peak within the first \textit{24~hours}
    \item Refrigerate in a sealed jar for up to \textit{4~days}
    \item Dressing will separate and thicken when cold; let sit \textit{10~minutes} at room temperature, then shake or whisk
    \item If flavor dulls after storage, refresh with 1~tsp. \textbf{raspberry shrub} or a squeeze of lemon before serving
\end{itemize}

\subsection*{Serving Suggestions}
\begin{itemize}
    \item Toss with arugula, frisee, radicchio, or mixed baby greens
    \item Pair with roasted beets, goat cheese, and toasted walnuts
    \item Spoon over roasted or smoked salmon with cucumber and herbs
    \item Use on a farro or lentil salad with fennel and red onion
\end{itemize}

\end{document}
